\documentclass[12pt,a4paper]{article}

\usepackage{fontspec}
\usepackage[russian]{babel}
\usepackage{geometry}
\usepackage{amsmath}
\usepackage{amssymb}
\usepackage{booktabs}
\usepackage{longtable}
\usepackage{array}
\usepackage{tabularx}
\usepackage{graphicx}
\usepackage{adjustbox}
\usepackage{float}
\usepackage{caption}
\usepackage{hyperref}
\usepackage{setspace}
\usepackage{ragged2e}

\setmainfont{Times New Roman}
\setsansfont{Arial}
\setmonofont{Courier New}

\geometry{margin=2.5cm}
\setlength{\parindent}{0cm}
\setlength{\parskip}{0.4em}
\renewcommand{\arraystretch}{1.2}
\setstretch{1.15}
\captionsetup{font=small,labelfont=bf}
\setlength{\emergencystretch}{2em}
\graphicspath{{report_figures/}}
\newcolumntype{Y}{>{\RaggedRight\arraybackslash}X}

\begin{document}

\begin{titlepage}
\thispagestyle{empty}
\setcounter{page}{1}
\begin{center}
{\large Федеральное государственное автономное образовательное учреждение высшего образования}\\[0.4cm]
{\Large НАЦИОНАЛЬНЫЙ ИССЛЕДОВАТЕЛЬСКИЙ УНИВЕРСИТЕТ}\\[0.25cm]
{\Large ``ВЫСШАЯ ШКОЛА ЭКОНОМИКИ''}\\[0.4cm]
{\large Факультет экономических наук}\\[0.2cm]
{\large Образовательная программа ``Экономика''}
\end{center}

\vspace{5.8cm}

\begin{center}
{\LARGE \textbf{Анализ динамических взаимосвязей между инфляцией, денежной массой и чистыми иностранными активами банковского сектора в России}}\\
\end{center}

\vfill

\begin{flushright}
\textbf{Выполнила:}\\[0.5cm]
Ковалёва Мария Сергеевна\\[0.5cm]
группа БЭК2310\\[2cm]

\textbf{Преподаватель семинаров:}\\
Захарова Светлана Владимировна
\end{flushright}

\vspace{1.5cm}

\begin{center}
Москва --- 2026
\end{center}
\end{titlepage}
\setcounter{page}{2}

\tableofcontents
\newpage

\section{Введение}
Инфляция, денежное предложение и состояние банковского сектора являются важнейшими
характеристиками макроэкономической среды. В экономической теории и прикладных
исследованиях большое внимание уделяется вопросу о том, каким образом изменение денежной
массы влияет на динамику цен и через какие каналы происходит передача денежных и финансовых
шоков в экономику. Для российской экономики данная проблема особенно актуальна, поскольку на
инфляционные процессы оказывают влияние не только внутренние монетарные факторы, но и
внешнеэкономические условия, валютный курс, международные потоки капитала и состояние
финансового сектора.

В период 2013--2015 гг. российская экономика столкнулась с рядом существенных внешних и
внутренних шоков: замедлением экономического роста, снижением мировых цен на нефть,
ослаблением курса рубля, изменениями денежно-кредитной политики Банка России и переходом к
режиму инфляционного таргетирования. В этих условиях особый интерес представляет
исследование взаимосвязей между инфляцией, денежной массой и внешней позицией банковского
сектора.

В работе используются месячные данные по Российской Федерации за период с января 2013 года
по июнь 2015 года. Анализ проводится для трех макроэкономических показателей: индекса
потребительских цен ($\mathrm{CPI}$), денежного агрегата $\mathrm{M2}$ и чистых иностранных активов кредитных
организаций ($\mathrm{NFA}$). Рассматриваемый набор показателей позволяет исследовать взаимосвязи
между ценовой динамикой, денежным предложением и внешним финансовым каналом функционирования
банковской системы.

\subsection{Научная статья}
В качестве основной научной статьи для исследования выбрана работа:

\begin{center}
\textbf{Zeynalov A., Tiron K. (2022). Macroeconomic performance of oil price shocks in
Russia. arXiv:2211.04954.}
\end{center}

В данной статье исследуется влияние внешних шоков, связанных с изменением мировых цен на
нефть, на макроэкономические показатели Российской Федерации. Авторы рассматривают динамику
промышленного производства, инфляции, обменного курса и процентных ставок, используя
квартальные данные за период с 2004 по 2021 год. Для анализа взаимосвязей между
показателями применяется модель векторной авторегрессии (VAR), позволяющая исследовать
реакцию макроэкономических переменных на различные экономические шоки и оценивать каналы их
распространения в экономике.

Результаты исследования показывают, что внешние шоки оказывают значимое влияние на
российскую экономику через валютный и денежно-кредитный каналы. Авторы приходят к выводу,
что динамика инфляции определяется не только внутренними факторами, но и изменениями
внешнеэкономической конъюнктуры, которые воздействуют на финансовую систему, процентные
ставки и обменный курс. Для анализа этих взаимосвязей используются инструменты, близкие к
применяемым в настоящей работе, включая многомерное моделирование временных рядов и
исследование динамических эффектов между экономическими показателями.

Выбор данной статьи в качестве основы исследования обусловлен несколькими причинами.
Объектом анализа выступает российская экономика, что обеспечивает сопоставимость
экономических условий и интерпретации результатов. В статье используется
VAR-подход, который также применяется в настоящем исследовании. Кроме того, работа
рассматривает механизмы передачи экономических шоков через финансовые и денежные каналы, что
концептуально соответствует задаче изучения взаимосвязей между инфляцией, денежной массой и
чистыми иностранными активами банковского сектора.

Дополнительно следует отметить работу \textbf{Albulescu C. (2020). Bank financial stability,
bank valuation and international oil prices: Evidence from listed Russian public banks}, в
которой показано, что внешние экономические шоки способны оказывать существенное влияние на
состояние банковского сектора России. Данная работа подтверждает важность учета показателей,
характеризующих внешнюю позицию банковской системы, при анализе макроэкономических
процессов.

\subsection{Проблема исследования}
Несмотря на большое количество исследований, посвященных инфляции и денежно-кредитной
политике, вопрос о характере взаимосвязей между денежной массой, инфляцией и внешней
позицией банковского сектора остается дискуссионным. Согласно монетаристскому подходу, рост
денежного предложения в долгосрочном периоде приводит к росту уровня цен. Однако для
современных экономик влияние денежной массы на инфляцию часто оказывается менее прямым и
может реализовываться через различные финансовые каналы.

Для российской экономики особый интерес представляет изучение роли банковского сектора в
процессе передачи денежных импульсов. Изменения объема денежной массы могут сопровождаться
изменениями внешней позиции кредитных организаций, отражаемой показателем чистых иностранных
активов. Одновременно инфляционные процессы могут определяться не только внутренними
денежными факторами, но и внешними экономическими шоками, изменениями валютного курса и
финансовых условий.

В связи с этим возникает исследовательский вопрос: существуют ли статистически значимые
краткосрочные и долгосрочные взаимосвязи между инфляцией, денежной массой и чистыми
иностранными активами кредитных организаций в России, а также каковы направления передачи
влияния между этими показателями?

Для ответа на данный вопрос в работе используются методы анализа многомерных временных
рядов, включая тесты на стационарность, анализ коинтеграции, тест причинности по Грейнджеру,
модели VAR и ECM, а также анализ функций импульсного отклика и разложения дисперсии ошибок
прогноза.

\subsection{Гипотеза исследования}
Основная гипотеза исследования заключается в том, что между индексом потребительских цен
($\mathrm{CPI}$), денежным агрегатом $\mathrm{M2}$ и чистыми иностранными активами кредитных организаций
($\mathrm{NFA}$) существуют статистически значимые взаимосвязи.

Предполагается, что увеличение денежной массы способствует изменению инфляционной динамики,
однако данный эффект может проявляться с временным лагом и быть относительно слабым в
краткосрочном периоде. Также предполагается наличие более тесной взаимосвязи между денежной
массой и чистыми иностранными активами банковского сектора, поскольку изменения внешней
позиции кредитных организаций влияют на ликвидность банковской системы и процессы денежного
обращения.

Таким образом, ожидается, что показатели $\mathrm{M2}$ и $\mathrm{NFA}$ будут демонстрировать более
выраженную взаимосвязь между собой, чем каждый из них с показателем $\mathrm{CPI}$. Кроме того,
предполагается существование направленного влияния со стороны денежной массы на показатели
финансового сектора и инфляции, которое может быть выявлено с помощью тестов причинности по
Грейнджеру и моделей векторной авторегрессии.

\subsection{Информация о расчетной выборке}
В расчетной части исследования используется месячный датасет, сформированный на основе трех
макроэкономических показателей по Российской Федерации. В итоговую таблицу вошли:
\begin{itemize}
\item индекс потребительских цен на все товары и услуги, база 2010 = 100;
\item чистые иностранные активы кредитных организаций, млрд руб.;
\item денежный агрегат $\mathrm{M2}$, млрд руб.
\end{itemize}

Выбор именно этих рядов для расчетной части обусловлен следующими соображениями.
Все три показателя относятся к одной макроэкономической системе и могут быть
содержательно связаны между собой: динамика денежной массы и состояние внешних активов
банковского сектора потенциально связаны с инфляционными процессами. Все ряды
имеют месячную частоту, что позволяет использовать многомерные модели временных рядов без
дополнительной агрегации. Все показатели заданы в уровнях, принимающих
положительные значения, поэтому к ним применимы логарифмические преобразования, которые
удобны для интерпретации темпов изменения и для перехода к стационарным формам.

Вся дальнейшая расчетная часть опирается на выборку с января 2013 года по июнь 2015 года,
то есть на 30 наблюдений. Исходные ряды были подготовлены в файлах
\texttt{raw\_data\_29.xls}, \texttt{raw\_data\_30.xls}, \texttt{raw\_data\_32.xls} и сведены
в итоговую таблицу \texttt{final\_dataset\_29\_30\_32.csv}.

\begin{table}[H]
\centering
\caption{Исходные данные расчетной части}
\begin{tabularx}{\textwidth}{>{\RaggedRight\arraybackslash}p{3.7cm} >{\RaggedRight\arraybackslash}p{2.3cm} >{\RaggedRight\arraybackslash}p{1.4cm} >{\RaggedRight\arraybackslash}p{2.0cm} Y}
\toprule
Показатель & Единица измерения & Страна & Период & Причина включения в расчетную часть \\
\midrule
Индекс потребительских цен на все товары и услуги (CPI) & индекс, 2010 = 100 & Россия & 2013-01 -- 2015-06 & Базовый индикатор инфляционной динамики; позволяет анализировать ценовые изменения в системе макроэкономических переменных \\
Чистые иностранные активы кредитных организаций (NFA) & млрд руб. & Россия & 2013-01 -- 2015-06 & Отражают состояние внешней позиции банковского сектора и могут быть связаны с денежным предложением и ценовой динамикой \\
Денежный агрегат M2 & млрд руб. & Россия & 2013-01 -- 2015-06 & Характеризует денежную массу в экономике и является естественным кандидатом для многомерных моделей денежно-инфляционных связей \\
\bottomrule
\end{tabularx}
\end{table}

\section{Описание данных}
На первом этапе была выполнена базовая техническая проверка набора данных:
корректность формата даты, отсутствие пропусков и проверка числовых типов. После загрузки
датасет содержит 30 строк и четыре столбца: дату и три анализируемых показателя. Такой этап
не является формальностью: многомерные модели особенно чувствительны к пропускам и к
неправильной упорядоченности наблюдений, поэтому до начала эконометрического анализа важно
исключить технические ошибки.

Затем были построены графики уровней рядов. По визуальной динамике можно сделать
несколько содержательных наблюдений. Ряд $\mathrm{CPI}$ демонстрирует устойчивый восходящий тренд:
индекс увеличивается почти монотонно, при этом ускорение роста особенно заметно в конце
выборки. Ряд $\mathrm{M2}$ также в целом возрастает, хотя внутри выборки присутствуют локальные
снижения и участки замедления. Ряд $\mathrm{NFA}$ ведет себя существенно менее гладко: в нем есть
резкие скачки, сильная амплитуда колебаний и заметная нестабильность масштаба изменений.
Уже на этой стадии становится ясно, что ряды в уровнях вряд ли можно считать стационарными,
а различия в характере динамики требуют перехода к преобразованным формам.

\begin{figure}[H]
\centering
\includegraphics[width=0.92\textwidth]{levels.png}
\caption{Динамика исходных рядов в уровнях, 2013-01--2015-06; ось $Y$: CPI --- индекс 2010 = 100, NFA и M2 --- млрд руб.}
\end{figure}

После графиков уровней были построены графики логарифмов рядов. Использование
логарифмов имеет две задачи. Первая --- сделать масштабы переменных сопоставимыми, поскольку
$\mathrm{M2}$ и $\mathrm{NFA}$ измеряются в миллиардах рублей, а $\mathrm{CPI}$ --- индекс. Вторая --- перейти к
приростным интерпретациям: разности логарифмов приближенно описывают темпы изменения
показателей. В логарифмах тренд у $\mathrm{CPI}$ и $\mathrm{M2}$ сохраняется, но визуально динамика становится
более ровной, а для $\mathrm{NFA}$ логарифмирование уменьшает масштаб колебаний, хотя не устраняет их
полностью.

\begin{figure}[H]
\centering
\includegraphics[width=0.92\textwidth]{logs.png}
\caption{Логарифмы исходных рядов, 2013-01--2015-06; ось $Y$: логарифм уровня ряда}
\end{figure}

Для описания масштаба переменных были рассчитаны основные статистики: среднее, стандартное
отклонение, минимум и максимум.

\begin{table}[H]
\centering
\caption{Основные описательные статистики исходных рядов, 2013-01--2015-06}
\begin{tabular}{lrrrr}
\toprule
Показатель & Среднее & Ст. откл. & Минимум & Максимум \\
\midrule
$\mathrm{CPI}$ & 130.8807 & 10.4907 & 118.4000 & 150.8600 \\
$\mathrm{NFA}$ & 2928.9272 & 1069.6285 & 1546.8940 & 5164.0190 \\
$\mathrm{M2}$ & 29590.0467 & 1580.3052 & 26348.7000 & 31958.0000 \\
\bottomrule
\end{tabular}
\end{table}

Из таблицы видно, что ряды сильно различаются по масштабу. Это еще один аргумент в пользу
работы с логарифмами и разностями, поскольку непосредственное сравнение уровней здесь
малоинформативно. Наиболее волатильным в абсолютном выражении выступает $\mathrm{NFA}$, что
согласуется с визуальным анализом.

Дополнительно были рассчитаны месячные темпы изменения по формуле
\[
g_t = \left(\frac{y_t}{y_{t-1}} - 1\right)\cdot 100.
\]
Для них были получены сводные характеристики.

\begin{table}[H]
\centering
\caption{Сводные характеристики месячных темпов изменения, 2013-01--2015-06, \%}
\begin{tabular}{lrrrr}
\toprule
Показатель & Среднее & Ст. откл. & Минимум & Максимум \\
\midrule
$\mathrm{CPI}$ & 0.8420 & 0.7956 & 0.1471 & 3.8484 \\
$\mathrm{M2}$ & 0.6893 & 2.1331 & -4.1546 & 7.9045 \\
$\mathrm{NFA}$ & 6.8177 & 37.4332 & -37.0450 & 191.8871 \\
\bottomrule
\end{tabular}
\end{table}

Эти результаты позволяют содержательно интерпретировать различия в динамике. Для $\mathrm{CPI}$
темпы изменения положительны почти на всем интервале и достаточно плавны, что соответствует
инфляционному процессу с ускорением в конце выборки. Для $\mathrm{M2}$ разброс намного больше:
денежная масса растет в среднем умеренно, но отдельные месяцы показывают как заметное
сокращение, так и сравнительно высокие положительные приросты. Для $\mathrm{NFA}$ среднее значение
темпа роста оказывается высоким, однако оно сопровождается чрезвычайно большой дисперсией;
это означает, что ряд содержит очень крупные всплески и падения, а значит, в уровнях он
наиболее проблематичен для стандартных линейных моделей.

Еще один этап --- анализ парных корреляций в уровнях.

\begin{table}[H]
\centering
\caption{Корреляционная матрица рядов в уровнях, 2013-01--2015-06}
\begin{tabular}{lrrr}
\toprule
 & $\mathrm{CPI}$ & $\mathrm{NFA}$ & $\mathrm{M2}$ \\
\midrule
$\mathrm{CPI}$ & 1.0000 & 0.6324 & 0.8632 \\
$\mathrm{NFA}$ & 0.6324 & 1.0000 & 0.6887 \\
$\mathrm{M2}$ & 0.8632 & 0.6887 & 1.0000 \\
\bottomrule
\end{tabular}
\end{table}

Наиболее сильная связь наблюдается между $\mathrm{CPI}$ и $\mathrm{M2}$; корреляция составляет примерно
$0.86$. При этом важно подчеркнуть, что такая корреляция еще не означает ни причинности, ни
долгосрочного равновесия. Для нестационарных рядов высокая корреляция часто отражает
совместный тренд, а не содержательную экономическую зависимость. Именно поэтому после
описательной части в работе последовательно выполняются тесты стационарности, коинтеграции и
причинности.

Для исходных уровней были построены автокорреляционные и частные автокорреляционные функции
для всех трех рядов в уровнях. Их качественная интерпретация следующая. Во всех случаях
автокорреляции на малых лагах остаются высокими, а затухание не выглядит быстрым. Для
макроэкономических рядов это типичный признак тренда или единичного корня. Следовательно,
уже по ACF и PACF можно ожидать, что перед оцениванием многомерных моделей потребуется
переход к разностям.

\begin{figure}[H]
\centering
\includegraphics[width=0.92\textwidth]{acf_pacf_1.png}
\caption{ACF и PACF для ряда CPI в уровнях, 2013-01--2015-06; ось $X$: лаги в месяцах}
\end{figure}

\begin{figure}[H]
\centering
\includegraphics[width=0.92\textwidth]{acf_pacf_2.png}
\caption{ACF и PACF для ряда M2 в уровнях, 2013-01--2015-06; ось $X$: лаги в месяцах}
\end{figure}

\begin{figure}[H]
\centering
\includegraphics[width=0.92\textwidth]{acf_pacf_3.png}
\caption{ACF и PACF для ряда NFA в уровнях, 2013-01--2015-06; ось $X$: лаги в месяцах}
\end{figure}

\subsection{Сезонность}
Поскольку данные имеют месячную частоту, потенциально в них могла бы присутствовать
сезонность. Визуально наиболее заметны не регулярные сезонные колебания, а тренд и
структурные изменения: CPI почти монотонно растет с ускорением в конце выборки, M2 имеет
общий восходящий тренд с отдельными скачками, а NFA характеризуется резкими нерегулярными
изменениями масштаба. Устойчивый повторяющийся годовой паттерн по графикам не выглядит
главной особенностью рядов.

Сезонная корректировка в работе не выполнялась. Причина состоит в короткой выборке: 30
месячных наблюдений дают всего около двух с половиной годовых циклов, поэтому анализ лага 12
и надежная оценка сезонных факторов были бы нестабильными. Поэтому дальнейшая расчетная
часть фокусируется на стационарных преобразованиях, краткосрочной динамике и кризисном
структурном эпизоде конца 2014 года. Это ограничение учитывается при интерпретации моделей.

\section{Стационарность}
Проверка стационарности выполнена с помощью трех тестов: ADF, Phillips--Perron (PP) и
KPSS. ADF и PP проверяют нулевую гипотезу о наличии единичного корня, тогда как KPSS,
наоборот, использует в качестве нулевой гипотезы стационарность ряда. Совместное
использование этих тестов позволяет не опираться на одну процедуру и лучше увидеть
пограничные случаи.

Для теста ADF оценивалась регрессия вида
\[
\Delta x_t = \alpha + \beta t + \rho x_{t-1} + \sum_{i=1}^{p}\phi_i \Delta x_{t-i} + \varepsilon_t,
\]
где проверяется гипотеза $H_0 \colon \rho = 0$ о наличии единичного корня. PP-тест имеет
ту же нулевую гипотезу о единичном корне, но использует непараметрическую поправку к
автокорреляции и гетероскедастичности ошибок.

Для каждого показателя проверялись три формы:
\begin{itemize}
\item логарифм уровня ряда, $\log y_t$;
\item первая разность логарифма, $\Delta \log y_t$;
\item вторая разность логарифма, $\Delta^2 \log y_t$.
\end{itemize}

Результаты сведены в таблицу в формате задания. Для ADF и PP запись ``$H_0$ откл.''
означает отвержение единичного корня; для KPSS та же запись означает отвержение
стационарности.

\begin{table}[H]
\centering
\caption{Проверка стационарности рядов, 2013-01--2015-06}
\scriptsize
\begin{adjustbox}{max width=\textwidth}
\begin{tabularx}{\textwidth}{llYYYY}
\toprule
Ряд & Преобразование & ADF: статистика, $p$, решение & PP: статистика, $p$, решение & KPSS: статистика, $p$, решение & Итоговый вывод \\
\midrule
CPI & $\log y_t$ & 3.0155; $p=1.0000$; $H_0$ не откл. & 1.8613; $p=0.9985$; $H_0$ не откл. & 0.8011; $p=0.0100$; $H_0$ откл. & Уровень нестационарен; CPI не рассматривается как надежный I(1)-ряд \\
CPI & $\Delta \log y_t$ & 2.4623; $p=0.9990$; $H_0$ не откл. & -1.5212; $p=0.5230$; $H_0$ не откл. & 0.2582; $p=0.1000$; $H_0$ не откл. & Неоднозначный результат: ADF/PP не отвергают единичный корень, KPSS не отвергает стационарность \\
CPI & $\Delta^2 \log y_t$ & -3.5098; $p=0.0077$; $H_0$ откл. & -3.7198; $p=0.0038$; $H_0$ откл. & 0.1032; $p=0.1000$; $H_0$ не откл. & Стационарная рабочая форма; интерпретируется как изменение месячной инфляции \\
NFA & $\log y_t$ & -2.0127; $p=0.2810$; $H_0$ не откл. & -1.7452; $p=0.4080$; $H_0$ не откл. & 0.6156; $p=0.0212$; $H_0$ откл. & Уровень нестационарен; ряд трактуется как I(1) \\
NFA & $\Delta \log y_t$ & -5.5420; $p=0.0000$; $H_0$ откл. & -7.6721; $p=0.0000$; $H_0$ откл. & 0.0475; $p=0.1000$; $H_0$ не откл. & Стационарная рабочая форма \\
NFA & $\Delta^2 \log y_t$ & -9.4412; $p=0.0000$; $H_0$ откл. & -20.8674; $p=0.0000$; $H_0$ откл. & 0.0381; $p=0.1000$; $H_0$ не откл. & Стационарно после избыточного дифференцирования; не основная форма \\
M2 & $\log y_t$ & -1.7227; $p=0.4194$; $H_0$ не откл. & -1.7740; $p=0.3934$; $H_0$ не откл. & 0.8121; $p=0.0100$; $H_0$ откл. & Уровень нестационарен; ряд трактуется как I(1) \\
M2 & $\Delta \log y_t$ & -6.9888; $p=0.0000$; $H_0$ откл. & -8.5441; $p=0.0000$; $H_0$ откл. & 0.2218; $p=0.1000$; $H_0$ не откл. & Стационарная рабочая форма \\
M2 & $\Delta^2 \log y_t$ & -3.3505; $p=0.0128$; $H_0$ откл. & -21.6425; $p=0.0000$; $H_0$ откл. & 0.1197; $p=0.1000$; $H_0$ не откл. & Стационарно после избыточного дифференцирования; не основная форма \\
\bottomrule
\end{tabularx}
\end{adjustbox}
\normalsize
\end{table}

Подробная интерпретация этих результатов важна для всей дальнейшей работы.

Для CPI ряд в логарифме уровня явно нестационарен: ADF и PP не отвергают наличие единичного
корня, а KPSS отвергает стационарность. Первая лог-разность оказывается пограничным случаем:
ADF и PP с константой не отвергают единичный корень, тогда как KPSS не отвергает
стационарность. Поэтому для дальнейшего моделирования выбран более осторожный вариант ---
вторая лог-разность $\Delta^2 \log \mathrm{CPI}_t$.

Важно подчеркнуть интерпретацию: $\Delta \log \mathrm{CPI}_t$ приближенно соответствует
месячной инфляции, а $\Delta^2 \log \mathrm{CPI}_t$ --- это уже изменение месячной инфляции,
или инфляционное ускорение. Следовательно, VAR в данной работе напрямую моделирует не саму
инфляцию, а изменение инфляционной динамики.

Для NFA и M2 вывод более однозначен: логарифмы уровней нестационарны, а первые лог-разности
стационарны по ADF, PP и KPSS. Поэтому эти два ряда рассматриваются как I(1), а в
краткосрочных моделях используются формы $\Delta \log \mathrm{NFA}_t$ и
$\Delta \log \mathrm{M2}_t$.

Итоговая рабочая таблица для многомерного моделирования имеет вид:
\[
\begin{aligned}
\text{d2log\_cpi}_t &= \Delta^2 \log(\mathrm{CPI}_t),\\
\text{dlog\_nfa}_t &= \Delta \log(\mathrm{NFA}_t),\\
\text{dlog\_m2}_t &= \Delta \log(\mathrm{M2}_t).
\end{aligned}
\]

После перехода к этим преобразованиям объем рабочей выборки снижается до 28 наблюдений,
поскольку разности уменьшают число доступных точек в начале ряда. Именно этот датасет затем
используется в тесте Грейнджера и в модели VAR.

\section{Наличие взаимосвязей}
\subsection{Коинтеграционная матрица}
Поскольку исходные ряды в уровнях нестационарны, следующим шагом становится проверка
долгосрочных связей между ними. Здесь важно методологическое ограничение: стандартные тесты
коинтеграции Энгла--Грейнджера и Йохансена корректнее применять к переменным одного порядка
интегрируемости, обычно I(1). По результатам тестов NFA и M2 трактуются как I(1), тогда как
для CPI рабочей формой в моделировании выбрана $\Delta^2 \log \mathrm{CPI}_t$. Поэтому
коинтеграционные проверки с участием CPI в уровнях сохраняются только как exploratory
результаты и не используются для сильных выводов о долгосрочном равновесии.

Коинтеграция анализируется двумя способами. Используется попарный тест Энгла--Грейнджера,
основанный на проверке стационарности остатков долгосрочной регрессии, а также многомерный
тест Йохансена для всей системы из трех переменных. Для полной системы результат также
трактуется как exploratory, поскольку CPI не рассматривается как надежный I(1)-ряд.

В попарной постановке оценивается уравнение
\[
y_t = \alpha + \beta x_t + u_t,
\]
после чего тестируется стационарность остатков $u_t$. В многомерной постановке используется
представление
\[
\Delta Y_t = \Pi Y_{t-1} + \sum_{i=1}^{k-1}\Gamma_i \Delta Y_{t-i} + u_t,
\]
где ранг матрицы $\Pi$ определяет число коинтеграционных соотношений в системе.

\begin{table}[H]
\centering
\caption{Попарные тесты Энгла--Грейнджера для логарифмов уровней, 2013-01--2015-06}
\begin{tabular}{llrrc}
\toprule
Левый ряд & Правый ряд & Статистика & $p$-value & Отвергается $H_0$ при 5\% \\
\midrule
CPI & NFA & -1.3674 & 0.8085 & Нет; exploratory \\
CPI & M2 & -2.0411 & 0.5071 & Нет; exploratory \\
NFA & M2 & -3.7742 & 0.0147 & Да; I(1) \\
\bottomrule
\end{tabular}
\end{table}

Попарные результаты показывают, что значимая коинтеграционная связь обнаруживается только
между NFA и M2. Именно эту пару можно интерпретировать содержательно, поскольку оба ряда
рассматриваются как I(1). Для пар с CPI гипотеза об отсутствии коинтеграции не отвергается,
но даже при другом результате вывод был бы только предварительным из-за проблемы порядка
интегрируемости CPI.

\begin{table}[H]
\centering
\caption{Коинтеграционная матрица для CPI, NFA и M2, 2013-01--2015-06}
\small
\begin{adjustbox}{max width=\textwidth}
\begin{tabularx}{\textwidth}{lYYY}
\toprule
 & CPI & NFA & M2 \\
\midrule
CPI & --- & Exploratory: EG $p=0.8085$, связи нет; CPI не трактуется как I(1) & Exploratory: EG $p=0.5071$, связи нет; CPI не трактуется как I(1) \\
NFA & Exploratory: EG $p=0.8085$, связи нет & --- & Да: EG $p=0.0147$; валидная попарная связь при I(1) \\
M2 & Exploratory: EG $p=0.5071$, связи нет & Да: EG $p=0.0147$; валидная попарная связь при I(1) & --- \\
\bottomrule
\end{tabularx}
\end{adjustbox}
\normalsize
\end{table}

Для проверки общей структуры системы используется тест Йохансена в спецификации
\texttt{det\_order = 0}, \texttt{k\_ar\_diff = 1}.

\begin{table}[H]
\centering
\caption{Тест Йохансена для системы CPI--NFA--M2, 2013-01--2015-06; exploratory}
\begin{tabular}{lrrrrrr}
\toprule
$H_0$ & Trace & Крит. 95\% & Reject & Max-eig & Крит. 95\% & Reject \\
\midrule
$r \le 0$ & 20.5325 & 29.7961 & Нет & 11.6739 & 21.1314 & Нет \\
$r \le 1$ & 8.8585 & 15.4943 & Нет & 8.6962 & 14.2639 & Нет \\
$r \le 2$ & 0.1623 & 3.8415 & Нет & 0.1623 & 3.8415 & Нет \\
\bottomrule
\end{tabular}
\end{table}

Многомерный тест не отвергает гипотезу $r = 0$, то есть коинтеграция для полной системы
трех переменных не подтверждается даже в exploratory-постановке. Следовательно, в текущей
расчетной постановке нет оснований строить VECM для всего набора рядов. Валидная попарная
связь между NFA и M2 позволяет рассматривать двухмерное ECM как дополнительную модель, но
долгосрочную коинтеграцию с участием CPI в работе не следует утверждать.

\subsection{Причинность по Грейнджеру}
Для стационарных преобразований рядов проверяется причинность по Грейнджеру. В качестве
основного результата используется лаг 1, поскольку именно VAR(1) выбрана как базовая
спецификация по BIC/HQIC и соображениям экономии степеней свободы. Проверка лагов 1--3
сохраняется ниже только как анализ чувствительности: перебор нескольких лагов и выбор
минимального $p$-value повышает риск data mining и множественного тестирования.

\begin{table}[H]
\centering
\caption{Основной тест причинности по Грейнджеру при лаге 1, 2013-03--2015-06}
\begin{tabular}{llrc}
\toprule
Причина & Следствие & Лаг & $p$-value \\
\midrule
$\Delta \log \mathrm{NFA}$ & $\Delta^2 \log \mathrm{CPI}$ & 1 & 0.8315 \\
$\Delta \log \mathrm{M2}$ & $\Delta^2 \log \mathrm{CPI}$ & 1 & 0.1133 \\
$\Delta^2 \log \mathrm{CPI}$ & $\Delta \log \mathrm{NFA}$ & 1 & 0.8646 \\
$\Delta \log \mathrm{M2}$ & $\Delta \log \mathrm{NFA}$ & 1 & 0.0002 \\
$\Delta^2 \log \mathrm{CPI}$ & $\Delta \log \mathrm{M2}$ & 1 & 0.9979 \\
$\Delta \log \mathrm{NFA}$ & $\Delta \log \mathrm{M2}$ & 1 & 0.8745 \\
\bottomrule
\end{tabular}
\end{table}

Единственная статистически значимая связь на уровне 5\% возникает по направлению
\[
\Delta \log \mathrm{M2}_{t-1} \longrightarrow \Delta \log \mathrm{NFA}_t.
\]
Это означает, что прошлые изменения денежной массы помогают улучшить прогноз текущих
изменений чистых иностранных активов, тогда как обратная направленность в данных не
подтверждается. Для CPI статистически надежной причинности не обнаружено: важно, что здесь
моделируется $\Delta^2 \log \mathrm{CPI}$, то есть изменение месячной инфляции, а не сама
инфляция.

\begin{table}[H]
\centering
\caption{Проверка чувствительности теста Грейнджера по лагам 1--3, 2013-03--2015-06}
\small
\begin{adjustbox}{max width=\textwidth}
\begin{tabular}{llrrr}
\toprule
Причина & Следствие & $p$(lag 1) & $p$(lag 2) & $p$(lag 3) \\
\midrule
$\Delta \log \mathrm{NFA}$ & $\Delta^2 \log \mathrm{CPI}$ & 0.8315 & 0.6272 & 0.7678 \\
$\Delta \log \mathrm{M2}$ & $\Delta^2 \log \mathrm{CPI}$ & 0.1133 & 0.0888 & 0.2343 \\
$\Delta^2 \log \mathrm{CPI}$ & $\Delta \log \mathrm{NFA}$ & 0.8646 & 0.9755 & 0.7936 \\
$\Delta \log \mathrm{M2}$ & $\Delta \log \mathrm{NFA}$ & 0.0002 & 0.0010 & 0.0042 \\
$\Delta^2 \log \mathrm{CPI}$ & $\Delta \log \mathrm{M2}$ & 0.9979 & 0.9653 & 0.9206 \\
$\Delta \log \mathrm{NFA}$ & $\Delta \log \mathrm{M2}$ & 0.8745 & 0.4355 & 0.2563 \\
\bottomrule
\end{tabular}
\end{adjustbox}
\normalsize
\end{table}

Таблица чувствительности показывает, что связь $\Delta \log \mathrm{M2} \to \Delta \log
\mathrm{NFA}$ сохраняется при лагах 1--3. При этом слабый результат для $\Delta \log
\mathrm{M2} \to \Delta^2 \log \mathrm{CPI}$ на лаге 2 ($p=0.0888$) не используется как
основной вывод, поскольку он незначим на 5\% и возникает только при переборе лагов.

С точки зрения дальнейшей спецификации моделей этот этап особенно важен. Если бы причинность
была двусторонней и одновременно наблюдалась коинтеграция всей системы, можно было бы
ожидать необходимость более сложной модели с тесной динамической взаимосвязью всех
переменных. В текущей работе структура оказывается намного проще: для краткосрочной
динамики существенен прежде всего канал от $\mathrm{M2}$ к $\mathrm{NFA}$, а для долгосрочной связи
дополнительно выделяется только пара $\mathrm{NFA}$--$\mathrm{M2}$.

\begin{figure}[H]
\centering
\includegraphics[width=0.72\textwidth]{granger_scheme.png}
\caption{Схема статистически значимых связей по тесту Грейнджера}
\end{figure}

Схема делает структуру взаимосвязей более наглядной: на уровне значимости 5\% остается
только один направленный канал, от изменений денежной массы к изменениям чистых
иностранных активов. Все остальные направления не показаны стрелками, поскольку
соответствующие $p$-value выше 0.05. Связь $M2 \to CPI$ можно рассматривать только как
слабый намек при более мягком уровне значимости, но не как надежный результат.

\section{Спецификации моделей}
На основании предыдущих этапов в расчетной части были выбраны несколько модельных
постановок.

Первая --- векторная авторегрессия VAR для системы
\[
y_t = \begin{pmatrix}
\Delta^2 \log \mathrm{CPI}_t \\
\Delta \log \mathrm{NFA}_t \\
\Delta \log \mathrm{M2}_t
\end{pmatrix}.
\]
Выбор VAR мотивирован тем, что после перехода к стационарным формам все три переменные можно
рассматривать как взаимозависимую систему краткосрочной динамики. VAR удобна тем, что не
требует заранее жестко задавать направление причинности: каждая переменная объясняется как
своими лагами, так и лагами остальных переменных.

Вторая --- VARX-спецификация с кризисной dummy-переменной для конца 2014 года. Она нужна
как проверка чувствительности базовой VAR к резкому валютно-финансовому эпизоду декабря
2014 --- января 2015 года.

Третья --- одноуравненная ADL/ARMAX-постановка для наиболее выраженного канала
$M2 \to NFA$. Она позволяет проверить, сохраняется ли результат VAR в более простой
динамической регрессии, где $\Delta \log \mathrm{NFA}_t$ объясняется собственным лагом и лагами
других переменных.

Четвертая --- двухмерная модель коррекции ошибок ECM для пары $\mathrm{NFA}$--$\mathrm{M2}$. Основание для нее
дает попарный тест Энгла--Грейнджера, который выявляет коинтеграцию именно этой пары.
Поскольку тест Йохансена не подтверждает коинтеграцию всей трехмерной системы, VECM для
трех переменных в расчетной части не оценивается.

Следует отдельно отметить ограничения используемых моделей.
\begin{itemize}
\item VAR описывает только краткосрочную динамику преобразованных стационарных рядов. Она
не дает непосредственного уравнения долгосрочного равновесия между всеми тремя уровнями
переменных.
\item ECM чувствительна к качеству коинтеграционного соотношения. Если долгосрочная связь
слаба или нестабильна, коэффициент при error-correction term может оказаться статистически
незначимым.
\item ADL/ARMAX является одноуравненной моделью и потому не описывает всю систему
одновременно. Она полезна как дополнительная проверка конкретного направления связи, но не
заменяет VAR.
\item Кризисная dummy-переменная выделяет всего два наблюдения, поэтому ее коэффициенты
следует трактовать как проверку чувствительности, а не как полноценную оценку структурного
режима.
\item VECM для всей системы не оценивалась, поскольку результаты теста Йохансена не дают
оснований считать трехмерную систему коинтегрированной.
\end{itemize}

\section{Моделирование}
\subsection{Выбор числа лагов в VAR}
Для VAR рассматривались лаги от 1 до 3. Предварительный анализ с помощью процедуры
\texttt{select\_order} показала, что разные критерии дают разные рекомендации: AIC тяготеет к
более длинной модели, BIC предпочитает более короткую, HQIC и FPE склоняются к лагу 1.
Это типичная ситуация для короткой выборки, когда увеличение числа лагов улучшает подгонку,
но одновременно быстро сокращает число степеней свободы.

\begin{table}[H]
\centering
\caption{Сравнение спецификаций VAR, 2013-03--2015-06}
\begin{tabular}{rrrrr}
\toprule
Лаг & AIC & BIC & HQIC & Устойчивость \\
\midrule
1 & -21.0375 & -20.4615 & -20.8662 & Да \\
2 & -20.7900 & -19.7738 & -20.4974 & Да \\
3 & -20.8326 & -19.3700 & -20.4269 & Да \\
\bottomrule
\end{tabular}
\end{table}

По совокупности критериев базовой спецификацией была выбрана VAR(1). У нее лучший
из рассматриваемых вариантов BIC и HQIC, а также наилучшее значение AIC среди моделей,
сравниваемых именно после оценки. Кроме того, при короткой выборке лаг 1 позволяет сохранить
достаточное число наблюдений для интерпретации коэффициентов.

\subsection{Оценка VAR(1)}
Оцененная модель имеет общий вид
\[
y_t = c + A_1 y_{t-1} + u_t,
\]
где $y_t = (\Delta^2 \log \mathrm{CPI}_t,\ \Delta \log \mathrm{NFA}_t,\ \Delta \log \mathrm{M2}_t)'$.

Численные оценки коэффициентов приведены ниже.

\begin{table}[H]
\centering
\caption{Коэффициенты VAR(1), 2013-03--2015-06}
\small
\begin{tabularx}{\textwidth}{>{\RaggedRight\arraybackslash}p{3.2cm} >{\centering\arraybackslash}p{3.2cm} >{\centering\arraybackslash}p{3.2cm} >{\centering\arraybackslash}p{3.2cm}}
\toprule
Переменная & Уравнение $\Delta^2 \log \mathrm{CPI}_t$ & Уравнение $\Delta \log \mathrm{NFA}_t$ & Уравнение $\Delta \log \mathrm{M2}_t$ \\
\midrule
const & -0.000832 & -0.029243 & 0.007798 \\
$L1.\Delta^2 \log \mathrm{CPI}$ & 0.237656 & -0.182211 & -0.017491 \\
$L1.\Delta \log \mathrm{NFA}$ & 0.004542 & 0.194659 & 0.003041 \\
$L1.\Delta \log \mathrm{M2}$ & 0.101666 & 8.159310 & -0.298575 \\
\bottomrule
\end{tabularx}
\normalsize
\end{table}

\begin{table}[H]
\centering
\caption{$p$-value коэффициентов VAR(1), 2013-03--2015-06}
\small
\begin{tabularx}{\textwidth}{>{\RaggedRight\arraybackslash}p{3.2cm} >{\centering\arraybackslash}p{3.2cm} >{\centering\arraybackslash}p{3.2cm} >{\centering\arraybackslash}p{3.2cm}}
\toprule
Переменная & Уравнение $\Delta^2 \log \mathrm{CPI}_t$ & Уравнение $\Delta \log \mathrm{NFA}_t$ & Уравнение $\Delta \log \mathrm{M2}_t$ \\
\midrule
const & 0.470065 & 0.459997 & 0.084820 \\
$L1.\Delta^2 \log \mathrm{CPI}$ & 0.209934 & 0.977676 & 0.981254 \\
$L1.\Delta \log \mathrm{NFA}$ & 0.354188 & 0.247701 & 0.874472 \\
$L1.\Delta \log \mathrm{M2}$ & 0.061618 & 0.000013 & 0.162210 \\
\bottomrule
\end{tabularx}
\normalsize
\end{table}

Содержательно главная особенность модели --- очень сильная и статистически значимая роль
лага $\Delta \log \mathrm{M2}$ в уравнении для $\Delta \log \mathrm{NFA}$. Коэффициент равен $8.1593$, а
$p$-value практически нулевой. Это согласуется с тестом Грейнджера и усиливает вывод о том,
что краткосрочные изменения денежной массы несут информацию для прогнозирования изменений
чистых иностранных активов. В остальных уравнениях статистически значимых коэффициентов на
уровне 5\% почти нет; наиболее близок к значимости лаг $\Delta \log \mathrm{M2}$ в уравнении для
$\Delta^2 \log \mathrm{CPI}$ с $p \approx 0.062$.

\subsection{VAR с кризисной dummy-переменной}
Для учета кризисного эпизода конца 2014 года была добавлена экзогенная переменная
\[
D_t =
\begin{cases}
1, & t = \text{декабрь 2014 или январь 2015},\\
0, & \text{иначе}.
\end{cases}
\]
Такая спецификация обозначается здесь как VARX(1), поскольку к базовой VAR(1) добавляется
экзогенный индикатор резкого макроэкономического шока.

\begin{table}[H]
\centering
\caption{Сравнение VAR(1) и VARX(1) с кризисной dummy, 2013-03--2015-06}
\begin{tabular}{lrrrrc}
\toprule
Модель & AIC & BIC & HQIC & Whiteness $p$ & Устойчивость \\
\midrule
VAR(1) & -21.0375 & -20.4615 & -20.8662 & 0.9065 & Да \\
VARX(1) + dummy & -21.2261 & -20.5062 & -21.0121 & 0.8819 & Да \\
\bottomrule
\end{tabular}
\end{table}

Добавление dummy немного улучшает информационные критерии AIC, BIC и HQIC, при этом модель
остается устойчивой, а тест на отсутствие автокорреляции остатков по-прежнему не выявляет
проблем. Нормальность остатков, как и в базовой модели, отвергается, поэтому этот блок
лучше трактовать как проверку чувствительности, а не как окончательную модель.

\begin{table}[H]
\centering
\caption{Коэффициенты кризисной dummy в VARX(1), 2013-03--2015-06}
\begin{tabular}{lrr}
\toprule
Уравнение & Коэффициент dummy & $p$-value \\
\midrule
$\Delta^2 \log \mathrm{CPI}_t$ & 0.0122 & 0.0047 \\
$\Delta \log \mathrm{NFA}_t$ & -0.2258 & 0.1760 \\
$\Delta \log \mathrm{M2}_t$ & 0.0241 & 0.2080 \\
\bottomrule
\end{tabular}
\end{table}

Наиболее заметный результат связан с уравнением для $\Delta^2 \log \mathrm{CPI}_t$: dummy
статистически значима и имеет положительный знак. Это согласуется с тем, что в конце 2014
--- начале 2015 года инфляционная динамика резко ускорилась. В уравнениях для $\mathrm{NFA}$ и $\mathrm{M2}$
коэффициенты dummy статистически незначимы на уровне 5\%. При этом ключевой лаг
$\Delta \log \mathrm{M2}$ в уравнении для $\Delta \log \mathrm{NFA}$ остается значимым, то есть основной вывод
о канале $M2 \to NFA$ не исчезает после контроля кризисного эпизода.

\subsection{Диагностика VAR(1)}
После оценки модели были выполнены три диагностических шага:
\begin{itemize}
\item проверка устойчивости;
\item тест на отсутствие автокорреляции остатков;
\item тест на нормальность распределения остатков.
\end{itemize}

Результаты диагностики:
\begin{itemize}
\item модель устойчива;
\item $p$-value теста whiteness равно 0.9065, поэтому гипотеза об отсутствии остаточной
автокорреляции не отвергается;
\item $p$-value теста нормальности округляется к нулю, следовательно, нормальность остатков
отвергается;
\item статистики Дарбина--Уотсона по уравнениям равны 1.6174, 2.3043 и 2.0930.
\end{itemize}

Такая картина интерпретируется следующим образом. С точки зрения динамической спецификации
VAR(1) выглядит приемлемо: она устойчива, а остатки не проявляют систематической
автокорреляции. Это означает, что лаговая структура модели, по крайней мере на базовом
уровне, захватывает основную краткосрочную зависимость в данных. Однако распределение ошибок
далеко от нормального, а значит, стандартную параметрическую интерпретацию нужно делать
осторожно. Для макроэкономической короткой выборки это не редкость, но такой результат
следует прямо фиксировать в отчете как ограничение качества модели.

\begin{figure}[H]
\centering
\includegraphics[width=0.9\textwidth]{var_residuals.png}
\caption{Остатки модели VAR(1), 2013-03--2015-06; ось $Y$: остатки}
\end{figure}

\subsection{Оценка ADL/ARMAX}
В качестве дополнительной одноуравненной проверки была оценена динамическая модель для
наиболее выраженного направления связи:
\[
\Delta \log \mathrm{NFA}_t =
\alpha + \phi \Delta \log \mathrm{NFA}_{t-1}
+ \beta_1 \Delta \log \mathrm{M2}_{t-1}
+ \beta_2 \Delta^2 \log \mathrm{CPI}_{t-1}
+ \delta D_t + u_t.
\]
Это ADL-спецификация в стационарных преобразованиях. Дополнительно была оценена
ARMAX-версия с авторегрессионной ошибочной компонентой порядка 1 и тем же набором
экзогенных переменных.

\begin{table}[H]
\centering
\caption{Сравнение ADL/ARMAX-спецификаций для $\Delta \log \mathrm{NFA}_t$, 2013-03--2015-06}
\small
\begin{tabular}{lrrrrrr}
\toprule
Модель & $R^2$ & Adj. $R^2$ & AIC & BIC & RMSE & DW \\
\midrule
ADL & 0.4576 & 0.3869 & -9.1883 & -4.0049 & 0.1760 & 2.3043 \\
ADL + dummy & 0.4993 & 0.4082 & -9.3458 & -2.8667 & 0.1691 & 2.2549 \\
ARMAX(1,0,0) + dummy & --- & --- & -3.5792 & 3.9694 & 0.1764 & 1.9667 \\
\bottomrule
\end{tabular}
\normalsize
\end{table}

\begin{table}[H]
\centering
\caption{Коэффициенты ADL с кризисной dummy, 2013-03--2015-06}
\begin{tabular}{lrr}
\toprule
Переменная & Коэффициент & $p$-value \\
\midrule
const & -0.0205 & 0.6080 \\
$\Delta \log \mathrm{NFA}_{t-1}$ & 0.2449 & 0.1630 \\
$\Delta \log \mathrm{M2}_{t-1}$ & 9.2158 & 0.0000 \\
$\Delta^2 \log \mathrm{CPI}_{t-1}$ & 4.0621 & 0.5740 \\
$D_t$ & -0.2258 & 0.1900 \\
\bottomrule
\end{tabular}
\end{table}

ADL подтверждает основной результат многомерного анализа: лаг $\Delta \log \mathrm{M2}$ остается
статистически значимым предиктором $\Delta \log \mathrm{NFA}$. Добавление crisis dummy немного
улучшает $R^2$, AIC и RMSE, однако сама dummy в ADL незначима на уровне 5\%. Поэтому ее
следует рассматривать как контроль на кризисный эпизод, а не как устойчивый отдельный
эффект. ARMAX дает тот же содержательный вывод по лагу $\mathrm{M2}$, но по AIC и BIC уступает
более простой ADL-модели, поэтому в качестве дополнительной одноуравненной спецификации
лучше оставить ADL с dummy.

\subsection{Оценка ECM}
Для пары $\mathrm{NFA}$--$\mathrm{M2}$ сначала оценивалось долгосрочное коинтеграционное
соотношение
\[
\log \mathrm{M2}_t = \alpha + \beta \log \mathrm{NFA}_t + \varepsilon_t.
\]
Получены оценки
\[
\widehat{\log \mathrm{M2}_t} = 9.432709 + 0.108762 \log \mathrm{NFA}_t,
\]
при этом коэффициент при $\log \mathrm{NFA}_t$ статистически значим, а коэффициент детерминации
долгосрочной регрессии составляет $R^2 = 0.552821$.

Затем на основе лагированных остатков $\widehat{\varepsilon}_{t-1}$ строится ECM:
\[
\Delta \log \mathrm{M2}_t = \gamma_0 + \gamma_1 \Delta \log \mathrm{NFA}_t + \gamma_2 \widehat{\varepsilon}_{t-1} + u_t.
\]

\begin{table}[H]
\centering
\caption{Результаты ECM для пары $\mathrm{NFA}$--$\mathrm{M2}$, 2013-02--2015-06}
\begin{tabular}{lrr}
\toprule
Коэффициент & Оценка & $p$-value \\
\midrule
const & 0.007347 & 0.0635 \\
$\Delta \log \mathrm{NFA}_t$ & -0.025470 & 0.2070 \\
$ECT_{t-1}$ & -0.088040 & 0.4998 \\
\bottomrule
\end{tabular}
\end{table}

Дополнительные характеристики модели: $R^2 = 0.154$, $Adj.\ R^2 = 0.089$,
$Prob(F) = 0.114$, $DW = 1.996$.

Эконометрическая интерпретация здесь сдержанная. Знак коэффициента при $ECT_{t-1}$ верный
для механизма корректировки: отрицательное значение означало бы возврат системы к
долгосрочному равновесию. Однако коэффициент статистически незначим, а общий объясняющий
потенциал модели невысок. Поэтому в расчетной части ECM разумно рассматривать как
дополнительную иллюстрацию возможной долгосрочной связи пары $\mathrm{NFA}$--$\mathrm{M2}$, но не как
основную модель исследования.

\subsection{Сводная таблица оцененных моделей}
Для выполнения требования о сравнении спецификаций все оцененные модели сведены в одну
таблицу. В ней указаны рабочие преобразования, ключевые эффекты, диагностика, информационные
критерии и ограничения интерпретации.

\begin{table}[H]
\centering
\caption{Сводная таблица оцененных моделей, 2013-01--2015-06}
\scriptsize
\renewcommand{\arraystretch}{1.08}
\setlength{\tabcolsep}{3pt}
\begin{adjustbox}{max width=\linewidth}
\begin{tabularx}{\linewidth}{p{2.1cm} Y Y Y Y}
\toprule
Модель & Переменные и спецификация & Ключевые эффекты & Диагностика и критерии & Вывод и ограничения \\
\midrule
VAR(1) &
$\Delta^2 \log \mathrm{CPI}$, $\Delta \log \mathrm{NFA}$, $\Delta \log \mathrm{M2}$; 1 лаг, без экзогенных переменных &
$L1.\Delta \log \mathrm{M2} \to \Delta \log \mathrm{NFA}$: coef $=8.1593$, $p=0.0000$ &
Устойчива; whiteness $p=0.9065$; нормальность отвергается. AIC $=-21.0375$, BIC $=-20.4615$ &
Основной краткосрочный канал: $\mathrm{M2} \to \mathrm{NFA}$. Ограничение: короткая выборка; $\Delta^2 \log \mathrm{CPI}$ означает ускорение инфляции \\
\addlinespace
VARX(1) + dummy &
Те же переменные; 1 лаг + dummy для 2014-12--2015-01 &
Dummy значима в уравнении $\Delta^2 \log \mathrm{CPI}$; канал $\mathrm{M2} \to \mathrm{NFA}$ сохраняется &
Устойчива; whiteness $p=0.8819$; нормальность отвергается. AIC $=-21.2261$, BIC $=-20.5062$ &
Кризисный эпизод важен для CPI-ускорения, но не отменяет результат по $\mathrm{M2} \to \mathrm{NFA}$. Ограничение: dummy отмечает только два наблюдения \\
\addlinespace
ADL/ARMAX &
Зависимая: $\Delta \log \mathrm{NFA}$; ADL(1) с лагами $\Delta \log \mathrm{M2}$, $\Delta^2 \log \mathrm{CPI}$, собственным лагом и dummy; ARMAX(1,0,0) как robustness &
В ADL $L1.\Delta \log \mathrm{M2}$: coef $=9.2158$, $p=0.0000$ &
ADL+dummy: DW $=2.2549$, JB $p=0.6824$, LB(4) $p=0.4489$; AIC $=-9.3458$, BIC $=-2.8667$. ARMAX: AIC $=-3.5792$, BIC $=3.9694$ &
Одноуравненная проверка поддерживает $\mathrm{M2} \to \mathrm{NFA}$. Ограничение: ADL не описывает всю систему; ARMAX хуже по критериям \\
\addlinespace
ECM &
$\log \mathrm{M2}$ и $\log \mathrm{NFA}$; попарная ECM для NFA--M2, уравнение для $\Delta \log \mathrm{M2}$ &
$ECT_{t-1}$: coef $=-0.0880$, $p=0.4998$ &
$R^2=0.154$, DW $=1.996$, Prob(F) $=0.114$; AIC $=-141.7$, BIC $=-137.6$ &
Механизм коррекции к равновесию статистически не подтвержден. Ограничение: коинтеграция только попарная; ECT незначим \\
\bottomrule
\end{tabularx}
\end{adjustbox}
\normalsize
\end{table}

\section{Интерпретация результатов}
\subsection{Интерпретация ECM}
Содержательно ECM показывает две вещи: краткосрочный эффект и эффект возврата к
равновесию. В текущей оценке краткосрочное влияние $\Delta \log \mathrm{NFA}_t$ на $\Delta \log \mathrm{M2}_t$
не подтверждается статистически, а коэффициент при $ECT_{t-1}$ тоже оказывается
незначимым. Это означает, что даже если долгосрочная связь между уровнями $\mathrm{NFA}$ и $\mathrm{M2}$
существует, в имеющейся короткой выборке скорость возврата к равновесию оценена слишком
неустойчиво, чтобы делать сильные выводы.

\subsection{Интерпретация VAR}
Для VAR(1) ключевая интерпретация связана с тем, что лаг изменений денежной массы
оказывается значимым предиктором изменений чистых иностранных активов. Иными словами,
динамика $\mathrm{M2}$ в прошлом месяце помогает прогнозировать динамику $\mathrm{NFA}$ в текущем месяце.
Обратный канал, а также влияние $\mathrm{CPI}$ на остальные ряды в текущей модели статистически
существенно не проявляются.

Такая конфигурация соответствует очень компактной структуре краткосрочных взаимодействий:
денежная масса выступает наиболее информативной переменной системы, тогда как CPI во второй
лог-разности играет скорее роль ряда, для которого сильных взаимных обратных связей на
данной выборке не найдено.

Проверка с кризисной dummy уточняет этот вывод. Индикатор декабря 2014 --- января 2015 года
значим в уравнении для $\Delta^2 \log \mathrm{CPI}_t$, что отражает резкое изменение инфляционной
динамики в кризисный эпизод. Однако в уравнении для $\Delta \log \mathrm{NFA}_t$ ключевой лаг
$\Delta \log \mathrm{M2}_{t-1}$ остается значимым и после добавления dummy. Следовательно, результат
по каналу $M2 \to NFA$ не является только следствием одного кризисного месяца.

\subsection{Интерпретация ADL/ARMAX}
ADL-модель служит дополнительной проверкой главного направления, найденного в VAR. В
одноуравненной постановке лаг $\Delta \log \mathrm{M2}_{t-1}$ также оказывается статистически
значимым в уравнении для $\Delta \log \mathrm{NFA}_t$. Это означает, что вывод о прогнозной роли
денежной массы для чистых иностранных активов сохраняется не только в системе VAR, но и в
более простой динамической регрессии.

Кризисная dummy в ADL имеет отрицательный знак, но не является статистически значимой на
уровне 5\%. Поэтому ее нельзя интерпретировать как надежный самостоятельный эффект на
$\mathrm{NFA}$; она скорее помогает проверить, не ломается ли основной коэффициент при $\mathrm{M2}$ после
учета конца 2014 года. ARMAX-спецификация дает похожий вывод по лагу $\mathrm{M2}$, но хуже по
информационным критериям, поэтому содержательно она выступает как проверка устойчивости, а
не как лучшая модель.

\subsection{Импульсные отклики}
Были построены функции импульсного отклика для горизонта 6 периодов без ортогонализации
шоков. Импульсный отклик показывает, как единичный шок в одной переменной распространяется
по системе в последующие месяцы. Формально отклик определяется как последовательность
матриц $\Psi_h$ в представлении
\[
y_t = \mu + \sum_{h=0}^{\infty}\Psi_h u_{t-h},
\]
где элемент матрицы $\Psi_h$ показывает реакцию одной переменной на шок другой через $h$
периодов.

Интерпретировать IRF следует осторожно. В данной работе используются неортогонализированные
импульсные отклики, поэтому они описывают динамику в рамках оцененной VAR(1), но не являются
структурной идентификацией причинных шоков. Результаты зависят от выбранной спецификации VAR
и короткой выборки; если использовать ортогонализированные шоки, интерпретация также зависела
бы от порядка переменных.

Наиболее интересен отклик $\Delta \log \mathrm{NFA}_t$ на шок в $\Delta \log \mathrm{M2}_t$, поскольку именно
эта связь уже была выявлена тестом Грейнджера и коэффициентами VAR. В качественном смысле
IRF согласуется с тем, что возмущение в денежной массе связано с последующей динамикой
чистых иностранных активов не только в момент оценки коэффициента, но и в нескольких
последующих периодах.
Для остальных связей отклики либо малы по масштабу, либо быстро затухают, что хорошо
согласуется с отсутствием статистически значимых коэффициентов в соответствующих уравнениях.

\begin{figure}[H]
\centering
\includegraphics[width=0.92\textwidth]{irf.png}
\caption{Функции импульсного отклика для модели VAR(1), горизонт 6 месяцев; ось $X$: горизонт, ось $Y$: отклик}
\end{figure}

\subsection{Разложение дисперсии ошибки прогноза}
Также было построено FEVD на горизонте 6 периодов. Это дает более формальную
интерпретацию относительной роли шоков разных переменных.

FEVD также следует читать как модельную декомпозицию внутри выбранной VAR-спецификации, а
не как доказательство структурной причинности. При другой спецификации лагов, другом наборе
переменных или ортогонализации шоков численные вклады могли бы измениться.

Для $\Delta^2 \log \mathrm{CPI}_t$ почти вся дисперсия ошибки прогноза остается собственной:
на горизонте 5 периодов около $86.3\%$ объясняется собственными шоками, а вклад шоков $\mathrm{M2}$
составляет примерно $13.2\%$. Следовательно, изменение месячной инфляции в используемой
преобразованной форме в основном определяется собственной внутренней динамикой, хотя вклад
денежной массы в небольшой степени просматривается.

Для $\Delta \log \mathrm{NFA}_t$ картина заметно интереснее. Уже на горизонте 1 периода собственные
шоки объясняют только около $54.5\%$ дисперсии ошибки прогноза, тогда как вклад шоков
$\Delta \log \mathrm{M2}_t$ достигает примерно $44.4\%$. К горизонту 5 этот вклад стабилизируется на
уровне около $44.7\%$. Это очень сильный результат в терминах FEVD: почти половина
неопределенности прогноза изменений $\mathrm{NFA}$ связана с инновациями в денежной массе.

Для $\Delta \log \mathrm{M2}_t$, наоборот, доминируют собственные шоки: на горизонте 5 их вклад
составляет около $92.1\%$, тогда как вклад $\mathrm{NFA}$ близок к $6.7\%$, а вклад $\mathrm{CPI}$ --- около
$1.2\%$. Иначе говоря, денежная масса в данной системе скорее влияет на другие переменные,
чем сама существенно объясняется ими.

В совокупности IRF и FEVD согласуются с главным выводом VAR-части: среди трех рядов наиболее
выраженный краткосрочный канал в выбранной спецификации проходит от динамики денежной массы к
динамике чистых иностранных активов.

\section{Заключение}
Проведенный анализ показал, что ряды $\mathrm{CPI}$, $\mathrm{NFA}$ и $\mathrm{M2}$ в уровнях не подходят для
непосредственного многомерного моделирования, поэтому дальнейшие расчеты выполнялись на
стационарных преобразованиях. Для $\mathrm{NFA}$ и $\mathrm{M2}$ рабочими формами стали первые лог-разности,
для $\mathrm{CPI}$ --- вторая лог-разность.

Попарный тест Энгла--Грейнджера выявил коинтеграцию только для пары $\mathrm{NFA}$--$\mathrm{M2}$, тогда как
тест Йохансена не подтвердил коинтеграцию всей трехмерной системы. Связи с CPI в
коинтеграционном блоке рассматриваются только как exploratory, поскольку рабочая форма CPI
предполагает вторую лог-разность, а стандартные коинтеграционные тесты рассчитаны прежде
всего на переменные одного порядка интегрируемости. Тест причинности по Грейнджеру при лаге
1 и модель VAR(1) показали наиболее выраженный краткосрочный канал от $\mathrm{M2}$ к $\mathrm{NFA}$.
Этот же вывод подтверждается одноуравненной ADL/ARMAX-проверкой: лаг изменения денежной
массы остается значимым предиктором изменения чистых иностранных активов. ECM для пары
$\mathrm{NFA}$--$\mathrm{M2}$ была оценена как дополнительная модель, однако статистически значимого механизма
корректировки к равновесию выявлено не было.

В целом полученные результаты позволяют рассматривать денежную массу как наиболее
информативную переменную системы, а чистые иностранные активы --- как показатель, наиболее
чувствительный к изменениям в денежном секторе. Прямая краткосрочная связь между денежной
массой и изменением месячной инфляции в рассматриваемой выборке выражена значительно слабее.
Добавление кризисной dummy для декабря 2014 --- января 2015 года показывает, что этот эпизод
важен для инфляционного ускорения, но не отменяет основного вывода о связи M2 $\to$ NFA.

\subsection{Сопоставление со статьей}
Сопоставление полученных результатов с работой Zeynalov A., Tiron K. (2022) следует
проводить аккуратно, поскольку совпадает не вся исследовательская постановка, а прежде всего
методологический каркас. В обоих исследованиях используется VAR-подход для анализа
макроэкономических взаимосвязей в российской экономике, однако состав переменных, временная
частота и исследуемый горизонт различаются. В статье рассматриваются квартальные данные за
длительный период 2004--2021 гг., причем в фокусе находятся нефтяные шоки, обменный курс,
процентные ставки и инфляция. В настоящей работе используется более короткий месячный
интервал 2013--2015 гг. и исследуется система из $\mathrm{CPI}$, $\mathrm{M2}$ и $\mathrm{NFA}$.

Несмотря на эти различия, качественное сопоставление возможно по двум направлениям.
Обе работы показывают, что инфляционные процессы не сводятся к простой
механической связи ``денежная масса $\to$ цены''. В статье подчеркивается значимость внешних
шоков и финансовых каналов передачи влияния на инфляцию. В настоящем исследовании также не
обнаружено устойчивой сильной краткосрочной причинности от денежной массы к CPI в выбранной
стационарной форме: коэффициенты в уравнении для $\Delta^2 \log \mathrm{CPI}_t$ в целом
статистически слабы, а основной тест Грейнджера на лаге 1 не дает значимой связи с изменением
месячной инфляции. Пограничное значение при лаге 2 ($p \approx 0.089$) рассматривается только
как robustness-наблюдение, а не как основной вывод.

Также обе работы подтверждают важность финансового канала. В статье он проявляется
через реакцию макроэкономических переменных на внешние шоки, которые передаются через
валютный и денежно-кредитный механизмы. В настоящей работе наиболее отчетливый результат
получен для связи $\Delta \log \mathrm{M2} \to \Delta \log \mathrm{NFA}$: она подтверждается и тестом
Грейнджера, и оценкой VAR(1), и ADL/ARMAX-проверкой, и разложением дисперсии ошибки
прогноза. Тем самым расчетная часть согласуется со статьей в том, что денежно-финансовые
взаимосвязи оказываются существеннее, чем простая прямая краткосрочная зависимость
инфляции от одной денежной переменной.

Одновременно различия результатов тоже закономерны. В статье центральным источником
возмущений выступают нефтяные и внешнеэкономические шоки, тогда как в данной работе такие
переменные не включены непосредственно в модель. По этой причине здесь не удается в полной
мере воспроизвести тот же механизм передачи влияния на инфляцию, который обсуждается в
статье. Кроме того, короткая выборка 2013--2015 гг. содержит период резких структурных
сдвигов в российской экономике, а это повышает чувствительность выводов к выбору лагов,
преобразований и спецификации тестов стационарности.

Следовательно, сопоставление со статьей носит не буквальный, а концептуальный характер.
Главное совпадение состоит в том, что для российской экономики денежные и финансовые каналы
действительно оказываются важными и не исчерпываются одной простой связью между денежной
массой и инфляцией. Главное различие состоит в том, что в настоящей сокращенной постановке
наиболее отчетливо идентифицируется канал от $\mathrm{M2}$ к $\mathrm{NFA}$, тогда как прямое статистически
устойчивое влияние на изменение месячной инфляции выражено существенно слабее.

\subsection{Сложности моделирования}
Главной трудностью расчетной части стала необходимость согласовать выбор стационарных
преобразований и одновременно сохранить экономическую интерпретацию рядов. Особенно это
касается $\mathrm{CPI}$, для которого первая лог-разность дала неоднозначный результат в зависимости
от включения тренда в тест ADF. Вторая трудность --- короткая выборка, ограничивающая число
лагов и усложняющая надежную оценку более насыщенных спецификаций. Наконец, дополнительной
проблемой стало то, что попарная коинтеграция не перенеслась на полную трехмерную систему,
из-за чего набор допустимых моделей пришлось сузить. Отдельно стоит отметить, что кризисная
dummy выделяет только два наблюдения, поэтому ее нельзя использовать для сильных выводов о
структурном режиме; она полезна именно как проверка чувствительности базовых результатов.

\section*{Список литературы}
\addcontentsline{toc}{section}{Список литературы}
\begin{enumerate}
\item Лекции и семинары курса ``Временные ряды и их практическое применение'', НИУ ВШЭ, 2026.
\item Zeynalov A., Tiron K. Macroeconomic Performance of Oil Price Shocks in Russia // arXiv preprint arXiv:2211.04954, 2022. Россия; VAR; инфляция, процентные ставки, курс рубля, промышленное производство. URL: \url{https://arxiv.org/abs/2211.04954}.
\item Albulescu C. Bank Financial Stability, Bank Valuation and International Oil Prices: Evidence from Listed Russian Public Banks // arXiv preprint arXiv:2004.12791, 2020. Россия; банковский сектор; внешние шоки и финансовая устойчивость банков. URL: \url{https://arxiv.org/abs/2004.12791}.
\item Drobyshevsky S., Trunin P. Inflation Determinants in Russia in 2000--2012.
\item Granville B., Mallick S. Does Inflation Targeting Matter for Output Growth? Evidence from Emerging Market Economies.
\item Korhonen I., Mehrotra A. Money Demand in Post-Crisis Russia.
\item Sims C. A. Macroeconomics and Reality // Econometrica. 1980. URL: \url{https://www.jstor.org/stable/1912017}.
\item Johansen S. Statistical Analysis of Cointegration Vectors. 1988.
\item Engle R., Granger C. Co-integration and Error Correction: Representation, Estimation and Testing. 1987.
\item Lane P., Milesi-Ferretti G. The External Wealth of Nations.
\item Lane P., Milesi-Ferretti G. International Financial Integration and Economic Growth.
\item Campa J., Goldberg L. Exchange Rate Pass-Through into Import Prices. 2002. URL: \url{https://www.nber.org/papers/w8934}.
\item Kilian L. Not All Oil Price Shocks Are Alike. 2009.
\end{enumerate}

\section*{Приложения}
\addcontentsline{toc}{section}{Приложения}
\subsection*{Приложение А. Файлы расчетной части}
В расчетной части использовались следующие файлы:
\begin{itemize}
\item \texttt{clear\_analysis.ipynb} --- основной воспроизводимый ноутбук с расчетами;
\item \texttt{final\_dataset\_29\_30\_32.csv} --- итоговый датасет;
\item \texttt{raw\_data\_29.xls}, \texttt{raw\_data\_30.xls}, \texttt{raw\_data\_32.xls} ---
исходные выгрузки рядов;
\item \texttt{read\_dataset.py} --- вспомогательный модуль загрузки данных.
\end{itemize}

\end{document}
